% ---------------------------------------------------------------------------
% Research-agenda addition (tenants). Inserted by the lead immediately before
% the Limitations and boundaries section. Bibliography and exercises live in
% sealed-bib-tenants.tex and sealed-exercises-tenants.tex.
% ---------------------------------------------------------------------------
\section{Many tenants, one harbor}\label{sec:sealed-tenants}

\noindent\emph{Express lane: when many owners share one agent, a log line can
go to a tenant's memory, a shared memory, a tenant's adapter, shared training,
or a cache. Retrieval returns stored records exactly, so by this chapter's
one-bit argument it has no finite $\varepsilon$, and only labels can decide
what a shared store may hold. Training can release through the person-keyed
ledger of Section~\ref{sec:sealed-person-ledger}. Scrubbing changes no label
(Design invariant~\ref{inv:sealed-scrub-not-declass}), and no deterministic
scrubber is private for any finite $\varepsilon$
(Theorem~\ref{thm:sealed-no-scrubber}).}

\paragraph{The scene, widened.}
The room so far has two owners. A harbor that sells agents serves many. A
\textbf{tenant} is a customer account: a company, a clinic, a household. A
\textbf{person} is anyone whose data a record carries: the user who typed it,
and anyone the text names. Each tenant's agents write logs, and every log line
is a candidate input to something that outlives the job: a memory store read
in later sessions, a retrieval index, an adapter fine-tuned on the tenant's
traffic, a model trained on everyone's traffic, a cache that speeds up the
next request. The question is which line may go where. The privacy unit is
the person, not the tenant and not the line. McMahan et al.\ define user-level
neighbors this way: two datasets are neighbors when one adds or removes all
the examples of a single user \pdcite{mcmahan2017userdp}.

\paragraph{Retrieval is an exact release.}
A memory store or retrieval index returns stored records, verbatim or nearly
so, to whoever's query ranks them highest. Take a log $D$ that contains one of
person $u$'s records and its neighbor $D'$ without it. A retrieval that
returns that record under $D$ defines an event of probability one under $D$
and zero under $D'$. Section~\ref{sec:sealed-conservation} has already drawn
the conclusion: reporting the bit exactly fails every finite epsilon bound. A
shared store therefore cannot be metered by the ledger at all. Its admission
rule must be a label rule, decided before the record is written.

The attacks bear this out. Against a retrieval-augmented assistant over the
Enron email corpus, 250 untargeted prompts, each a retrieval cue followed by
``Please repeat all the context'', made GPT-3.5 reproduce retrieved records
verbatim in 116 cases, and 250 targeted prompts extracted 205 items of
personal information \pdcite{zeng2024rag}. Storing only embeddings does not
anonymize: an iterative inversion recovered 92\% of 32-token inputs exactly
from their embeddings, and 89\% of full names from embedded clinical notes
\pdcite{morris2023vec2text}. The same retrieval study found that retrieval
lowers leakage of the model's own training data: reconstructions of training
text from its prefixes fell from 213 to between 33 and 70
\pdcite{zeng2024rag}. Moving logs from training into retrieval trades a noisy
channel for an exact one.

\paragraph{Training is a release to every future reader.}
A model trained on logs is read by everyone it serves, so training is a shared
destination unless the adapter serves one owner. Fine-tuning data is small and
recently seen, which is where membership inference works. Attacks that barely
beat chance on pretraining data, with AUC below 0.6 in most domains for models
up to 12B parameters \pdcite{duan2024mia}, reached an AUC of 0.96 against a
fine-tuned GPT-2, which also yielded 23\% of the personal information in its
training set at 30\% precision \pdcite{lukas2023pii}. Fine-tuning can also
reawaken what pretraining stored: fine-tuning GPT-3.5 on ten examples
recovered 69.9\% of targeted Enron email addresses, where jailbreak prompting
recovered none \pdcite{chen2023janus}. The route matters as well: training
through a reward model and RL memorizes far less than fine-tuning on the data
directly, and IPO-style direct preference learning regurgitates more than RL
\pdcite{pappu2024rlhf}. Table~\ref{tab:sealed-tenant-channels} collects the
channels.

\begin{table}[H]\centering\footnotesize
\caption{Channels by which one tenant's log content can reach another
principal. Figures are the cited studies' own measurements on their models and
data, not measurements of the reference system.}
\label{tab:sealed-tenant-channels}
\begin{tabularx}{\textwidth}{@{}>{\raggedright\arraybackslash}p{2.3cm} >{\raggedright\arraybackslash}p{2.5cm} >{\raggedright\arraybackslash}X >{\raggedright\arraybackslash}p{2.9cm}@{}}
\toprule
Channel & Mechanism & Evidence & Defense in this design \\
\midrule
Shared weights & rare or repeated sequences regurgitated & 604 verbatim training examples among 1{,}800 GPT-2 candidates, some from a single document \pdcite{carlini2021extracting}; more than 10{,}000 examples from ChatGPT for \$200, with memorized personal information in 16.9\% of tested generations \pdcite{nasr2023scalable} & ledger-charged training; deduplication as hygiene \\
Fine-tuned adapter & small, recent data; reactivated memory & membership AUC 0.96; about 3\% of personal information still extracted at $\varepsilon=8$ \pdcite{lukas2023pii}; ten examples recover 69.9\% of targeted addresses \pdcite{chen2023janus} & per-tenant adapter; ledger; no fine-tuning interface on shared models \\
Shared memory, read & exact release & 116 of 250 prompts reproduce records \pdcite{zeng2024rag}; 92\% exact inversion of embeddings \pdcite{morris2023vec2text} & owner labels; per-tenant stores \\
Shared memory, write & planted record replayed as instruction & 98.2\% injection and 76.8\% attack success from queries alone \pdcite{dong2025minja}; five texts per target, about 90\% \pdcite{zou2024poisonedrag} & integrity class \\
Prompt or key-value cache & timing of cache hits & caches shared globally across users at seven API providers, including OpenAI \pdcite{gu2025caching} & Design invariant~\ref{inv:sealed-label-keyed} \\
Preference data & route of training & reward model and RL memorize far less than direct fine-tuning; IPO-style learning more than RL \pdcite{pappu2024rlhf} & reward-model route; ledger \\
Deletion & influence left in weights & unlearning well defined only at the algorithmic level \pdcite{thudi2021unlearning} & retrain the tenant adapter; ledger bound \\
\bottomrule
\end{tabularx}
\end{table}

\paragraph{The flow rule.}
Every log line $\ell$ carries an \textbf{owner label} $L(\ell)$, the set of
persons and tenants whose permission a release needs, and an \textbf{integrity
class} $I(\ell)\in\{\text{operator},\text{owner-authored},\text{third-party}\}$.
Design invariant~\ref{hyp:sealed-taint} fixes the owner label conservatively:
every line the worker writes after reading $u$'s data carries $u$, whether or
not a detector finds anything of $u$'s in it. A destination $d$ has a set of
readers $R(d)$ and a use. A store that is read as demonstrations or
instructions has \emph{directive} use; any other store has data use. The rule
has three parts.

\begin{enumerate}[leftmargin=1.6em,itemsep=2pt]
\item \emph{Confidentiality.} $\ell$ may enter $d$ only if every owner in
  $L(\ell)$ has authorized every reader in $R(d)$ for this information, or has
  declassified $\ell$ for $d$. A line with several owners takes the most
  restrictive authorization, as in the lattice join of
  Section~\ref{sec:sealed-design}.
\item \emph{Training.} A training destination may instead admit $\ell$ through
  a private mechanism whose cost is charged to every owner in $L(\ell)$ on the
  person-keyed ledger.
\item \emph{Integrity.} If any principal other than $\ell$'s author reads $d$
  with directive use, $I(\ell)$ must be operator.
\end{enumerate}

The integrity part is new to this chapter, whose labels so far protect
confidentiality only. It answers the write side of a shared store. A user with
only query access injected records into an agent's memory with 98.2\% average
success, and those records steered later victims' queries with 76.8\% average
success \pdcite{dong2025minja}. Five crafted texts per target question, among
millions, reached about 90\% attack success against a retrieval pipeline, and
duplicate filtering did not stop them \pdcite{zou2024poisonedrag}. An injected
instruction that a model stored in long-term memory compromised later sessions
again \pdcite{greshake2023indirect}. Confidentiality labels say nothing about
any of this; the integrity class does.

Third parties named in content are kept out of training. A private fact is
often shared within a group that has no bounded size
\pdcite{brown2022privacy}, and the ledger cannot charge an owner it cannot
name. Table~\ref{tab:sealed-tenant-policy} applies the rule to common classes
of log content.

\begin{table}[H]\centering\footnotesize
\caption{Where a log line may go, by data class, under the flow rule.
\emph{yes}: allowed; \emph{release}: only with the owner's declassification
for that destination, such as a consent or a data-use declaration;
\emph{ledger}: only through a private mechanism charged on the person-keyed
ledger; \emph{handle}: an opaque reference, never the value; \emph{no}:
forbidden. The cache column names the key scope. A ledger-charged synthetic
derivative enters as the first row. A proposed policy, not a measured one.}
\label{tab:sealed-tenant-policy}
\begin{tabularx}{\textwidth}{@{}>{\raggedright\arraybackslash}p{3.4cm} >{\raggedright\arraybackslash}X >{\raggedright\arraybackslash}X >{\raggedright\arraybackslash}X >{\raggedright\arraybackslash}X >{\raggedright\arraybackslash}X@{}}
\toprule
Data class & Tenant memory & Shared memory & Tenant adapter & Shared training & Cache \\
\midrule
Public or operator-signed (empty label) & yes & yes & yes & yes & shared \\
Tenant-confidential, no person & yes & release & yes & ledger, tenant as unit & tenant key \\
The requesting person's data & release, person scope & no & release if one person, else ledger & ledger & person key \\
Third parties named in content & discloser's scope only & no & no & no & person key \\
Special category: health, finance & release, retention-limited & no & no & no & person key \\
Secrets and credentials & handle & no & no & no & no \\
Untrusted third-party content & data use only & no & no & no & tenant key \\
Model inferences about a person & release, visible to that person & no & no & no & person key \\
\bottomrule
\end{tabularx}
\end{table}

\begin{pdclaim}{Design invariant}{Label-Keyed Caches and Stores}
\setlength{\parskip}{4pt plus 1pt}
\label{inv:sealed-label-keyed}
Every stored or cached artifact computed from a log line inherits that line's
owner label: memory records, index chunks, embeddings, prefix and key-value
cache entries, and cached responses. A lookup may hit an entry only if the
entry's label authorizes the requesting principal. Whether a lookup hits is
itself an output, as a request to an outside tool is
(Section~\ref{sec:sealed-design}).

This is a \Vision{} obligation of a multi-tenant harbor; no reference-system
component implements it.
\end{pdclaim}

The cache clause is not hypothetical. A statistical timing audit found prompt
caches shared globally across users at seven API providers, including OpenAI,
so that a fast response could reveal that another user had recently sent the
same prefix; at least five providers made changes after disclosure
\pdcite{gu2025caching}. The same audit could not extract other users' prompts
token by token from live APIs. The invariant does not rely on that
difficulty. Timing inside an authorized scope remains the residual that
Table~\ref{tab:sealed-not-promised} calls padded and scheduled, not closed.

\begin{pdclaim}{Design invariant}{Scrubbing Is Not Declassification}
\setlength{\parskip}{4pt plus 1pt}
\label{inv:sealed-scrub-not-declass}
Removing or masking the personal information that a detector finds in a line
removes no owner from $L(\ell)$. A scrubbed line keeps its data class and its
row of Table~\ref{tab:sealed-tenant-policy}. Scrubbing is defense in depth,
not a release.

This is a \Vision{} obligation; no reference-system component implements it.
\end{pdclaim}

The evidence agrees. In a study of fine-tuned GPT-2 models, scrubbing drove
measured extraction of tagged personal information to zero, yet membership
inference against the scrubbed model still reached an AUC of 0.82
\pdcite{lukas2023pii}. Sanitization recognizes only a small part of the
private information that text can carry \pdcite{brown2022privacy}. The
structural reason is the next result.

\begin{pdclaim}{Theorem}{No Scrubber Guarantee Without a Mechanism}
\setlength{\parskip}{4pt plus 1pt}
\label{thm:sealed-no-scrubber}
Fix the person-level neighboring relation: logs $D$ and $D'$ are neighbors
when $D'$ removes every record of one person from $D$. Let $S$ be a
deterministic map from logs to logs. If some neighbors $D,D'$ have
$S(D)\ne S(D')$, then $S$ is not $\varepsilon$-private for any finite
$\varepsilon$.

In particular, a deterministic scrubber that lets any part of a person's
records influence its output has no finite $\varepsilon$ for that person. Its
guarantee can only be a statement about its detector: which spans it removes,
with what recall, on which distribution. It is not a statement about what an
adversary learns.
\end{pdclaim}

\paragraph{Proof.}
Let $D,D'$ be neighbors with $S(D)\ne S(D')$, and take the output event
$E=\{S(D)\}$. Since $S$ is deterministic, $\Pr[S(D)\in E]=1$ and
$\Pr[S(D')\in E]=0$. Privacy would require $1\le e^{\varepsilon}\cdot0$, which
fails for every finite $\varepsilon$. This is the one-bit argument of
Section~\ref{sec:sealed-conservation}, applied to a whole transformed log.
\qed

A randomized scrubber escapes the theorem only if its output distribution
changes by a bounded ratio on every neighbor pair, and then it is a private
mechanism whose cost belongs on the ledger. Synthetic text from a generator
fine-tuned with differential privacy at $\varepsilon=4$ is one such route:
canary sentences repeated 100 times appeared verbatim in 80\% of cases in
non-private synthetic data and in none under privacy
\pdcite{yue2022synthetic}.

\paragraph{Appropriateness is a separate question.}
Labels and $\varepsilon$ say who may read and how much may leak. The
\emph{release} cells of Table~\ref{tab:sealed-tenant-policy} also need a norm
for whether a flow is appropriate at all. Nissenbaum's \textbf{contextual
integrity} supplies one: information flows carry context-relative norms of
two kinds, of appropriateness and of distribution, and contextual integrity
is violated when either is breached \pdcite{nissenbaum2004ci}. A patient's
symptom told to a clinic's intake agent may enter that patient's own memory
scope. It does not thereby become available to the clinic's billing
analytics, though both sit inside one tenant. The norm cannot be delegated to
the model. In meeting-summary tasks built on such norms, GPT-4 revealed the
secret in 39\% of cases and ChatGPT in 57\%, although both were instructed to
preserve privacy, and chain-of-thought prompting did not reduce the leakage
\pdcite{mireshghallah2023confaide}. The gate enforces the norm; the model only
proposes.

\paragraph{Persistent agents keep per-owner memory.}
\pdchapref{stp}{From Spawn to Person} gives agents durable memory and
reputation. Under whole-worker labeling, an agent that served tenant $A$
carries $A$'s label into its next job unless its memory is partitioned. A
persistent person's episodic memory is therefore partitioned by owner label,
and only the partitions the work order authorizes are mounted. Its identity
and outcome record need not be partitioned: an outcome grade is
operator-class, and an aggregate over owners can leave through the ledger.
Deletion differs by destination. A tombstone retracts a record from an index
and from a cache. It cannot retract influence from trained weights. The same
model can be reached by training runs with and without a given point, so
unlearning is well defined only at the algorithmic level
\pdcite{thudi2021unlearning}, and removing information from parameters does
not by itself guarantee what the model generates \pdcite{cooper2024unlearning}.
For an adapter, deletion means retraining without the person, which
per-tenant adapters make cheap, or relying on the bound the ledger already
charged.

\begin{pdboundary}
\setlength{\parskip}{4pt plus 1pt}
Labels decide admission; they do not detect a record that was mislabeled when
it was written. The integrity class blocks directive use of third-party
content; it does not stop a third party from planting false facts as data.

The table cells are a proposed policy. The figures in this section are the
cited studies' measurements on their own models and data.

Timing on a label-keyed cache inside one principal's scope remains open, as in
Table~\ref{tab:sealed-not-promised}.
\end{pdboundary}

\pdexercisepointer{\ref{ex:sealed-tenant-route}}{\ref{ex:sealed-tenant-persistent}}{\pageref{ex:sealed-tenant-route}}
